
\chapter{Technical Analysis of \texttt{linkdraw.m}}
\date{March 10, 2026}


\maketitle

\section{Overview}
The \texttt{linkdraw.m} file visualizes directed graphs in 2D or 3D space based on a linking map $\phi$[cite: 4]. It performs geometric visualization by drawing vectors from start positions to their mapped positions[cite: 4].

\section{Step-by-Step Analysis}

\subsection{Section 1: Coordinate Standardization}
The function standardizes the input \texttt{pos}, supporting complex vectors (2D) or $N \times 3$ matrices (3D)[cite: 4].
\begin{lstlisting}
if ncol==2 
    pos=pos(:,1)+1i*pos(:,2);
elseif ncol>=3 
    drawdim=3; 
end
\end{lstlisting}

\subsection{Section 2: Rendering}
The code utilizes \texttt{quiver} or \texttt{quiver3} to render the graph edges and \texttt{text} to place labels at specified coordinates[cite: 4].
\begin{lstlisting}
hq=quiver(real(pos),imag(pos),real(dir),imag(dir),0);
ht=text(real(place),imag(place),num2str(labels));
\end{lstlisting}

\section{Expected Outputs and Visuals}
The code generates a Quiver Plot[cite: 4]. Edges represent relationships where an arrow points from $pos(i)$ to $pos(\phi(i))$[cite: 4].
